\documentclass[conference]{IEEEtran}

\usepackage{iftex}
\ifPDFTeX
\usepackage[utf8]{inputenc}
\usepackage[T1]{fontenc}
\fi
\usepackage{amsthm}
\usepackage{amsmath,amssymb,bm,mathtools}
\usepackage{newtxtext,newtxmath}
\usepackage{textcomp}
\usepackage{array,booktabs}
\usepackage{siunitx}
\usepackage{cite}
\usepackage[breaklinks=true,hidelinks]{hyperref}
\usepackage{bookmark}
\usepackage[nameinlink,capitalise,noabbrev]{cleveref}
\usepackage[final]{microtype}

\interdisplaylinepenalty=2500
\allowdisplaybreaks
\emergencystretch=2em
\sisetup{detect-weight=true,detect-family=true,per-mode=symbol}

\DeclareRobustCommand{\vct}[1]{\bm{#1}}
\DeclareRobustCommand{\mat}[1]{\bm{#1}}
\newcommand{\R}{\mathbb{R}}
\newcommand{\Id}{\mat I}
\newcommand{\norm}[1]{\left\lVert#1\right\rVert}
\newcommand{\abs}[1]{\left|#1\right|}
\newcommand{\T}{^{\mathsf T}}
\DeclareMathOperator{\blkdiag}{blkdiag}
\DeclareMathOperator{\diag}{diag}
\DeclareMathOperator{\Exp}{Exp}

\newtheorem{theorem}{Theorem}
\newtheorem{lemma}{Lemma}
\newtheorem{corollary}{Corollary}
\newtheorem{proposition}{Proposition}
\theoremstyle{remark}
\newtheorem{remark}{Remark}
\newtheorem{assumption}{Assumption}

\title{Pseudo-Measurement Anchoring and Input-to-State Stability in Marine Inertial Observers\\
\large A Unified Analysis of OU--II, OU--III, Two-Frame EKF, PII, and Time-Varying-Gain Nonlinear Estimators (draft)}
\author{%
  \IEEEauthorblockN{Mikhail S. Grushinskiy}
  \IEEEauthorblockA{Independent Researcher, USA, Fair Lawn\\
  Bareboat Necessities GitHub Project\\
  Email: mgrouch@users.github.com}%
}

\begin{document}
\maketitle

\begin{abstract}
Marine inertial motion estimators commonly suppress low-frequency integration
drift by introducing information that is not a conventional instantaneous
sensor measurement: zero-mean displacement constraints, integral-displacement
pseudo-measurements, virtual integrated-height measurements, or deterministic
restoring feedback.  This paper gives a common stability interpretation for
five estimators implemented in the same marine IMU code base: OU--II and
OU--III quaternion error-state Kalman filters, a right-invariant two-frame
Lie-group EKF (TFG), a vertical PII observer, and a nonlinear time-varying-gain
observer derived from the marine INS construction of Bryne, Fossen and
Johansen.  The central result is that bounded-gap anchoring of the head of an
integrator chain creates a uniformly observable finite-horizon block; stable
unobserved tails preserve detectability; and bounded exogenous scheduling plus
standard Riccati hypotheses yields uniform exponential stability of the
linearized Kalman recursion.  A quadratic nonlinear remainder then gives a
local input-to-state stability (ISS) bound.  The same structure survives smooth
bounded coordinate changes, which transfers the OU--III local result to the
TFG linearization under a Jacobian-consistency hypothesis.  For the PII
observer a scaled-coordinate Lyapunov argument gives an explicit slow-scheduling
condition.  Finally, the published time-varying-gain nonlinear-observer proof is
recast as an interconnection of two ISS subsystems and extended to expose the
additional schedule-rate term created when the high-gain scale itself is made
time varying.  The analysis also separates stability from accuracy: a virtual
zero constraint that is only approximately true appears as an input in the ISS
bound, so stronger anchoring can improve drift stability while simultaneously
biasing the desired wave-band signal.
\end{abstract}

\begin{IEEEkeywords}
input-to-state stability, pseudo-measurement, inertial navigation, marine
heave, error-state Kalman filter, invariant EKF, nonlinear observer, virtual
measurement, Ornstein--Uhlenbeck process
\end{IEEEkeywords}

\section{Introduction}

A free inertial integrator has neutral modes.  In a marine wave estimator this
is not merely a numerical inconvenience: accelerometer bias, attitude error and
low-frequency specific-force leakage are integrated into velocity and
position, eventually overwhelming the oscillatory motion of interest.  Several
practical estimators in this repository attack the same problem with apparently
different mechanisms:
\begin{itemize}
  \item OU--II applies periodic $p=0$ and $v=0$ pseudo-measurements to a
        $[v,p,a_w]$ translational chain;
  \item OU--III adds $S=\int p\,dt$ and applies the higher-order constraint
        $S=0$ to $[v,p,S,a_w]$;
  \item the two-frame group EKF retains the OU--III $S=0$ information in
        right-invariant coordinates;
  \item the vertical PII observer continuously feeds back $S,p,v$ with a
        repeated-real-pole restoring law;
  \item the nonlinear observer of Bryne, Fossen and Johansen augments vertical
        position with its integral and uses the virtual measurement of that
        integral at zero \cite{BryneFossenJohansen2014_TimeVarGain}.
\end{itemize}

The common object is therefore not the OU process, the Kalman filter, or a
particular attitude representation.  It is \emph{integral-state anchoring}:
information is injected at one or more levels of an otherwise drifting
kinematic chain.  Once that information is persistent enough to make the chain
uniformly observable or detectable, the remaining stability argument depends
on the estimator class rather than on the marine interpretation of the
constraint.

This paper has four goals.  First, it states a general finite-horizon
observability result for an integrated chain with either an OU tail or the
pure-integrator limit.  Second, it places the existing OU--II and OU--III local
ISS proofs under one Kalman/LTV theorem and shows which assumptions are shared.
Third, it treats TFG and PII without pretending they are ordinary instances of
the same Riccati recursion: TFG is handled by bounded local coordinate
conjugacy, while PII is handled by a direct Lyapunov proof.  Fourth, it gives an
ISS reading of the time-varying-gain nonlinear observer and identifies exactly
what changes when the repository schedules the high-gain scale $\theta$ rather
than only the multiplier $\vartheta(t)$ used as time varying in the published
proof.

The main conceptual limitation is equally important.  Stability of an estimator
that is anchored by the virtual equation ``integral displacement equals zero''
does not prove that the physical integral is zero at every instant.  Any
mismatch between the virtual equation and the actual wave/current/tide motion
enters the estimation error as a disturbance.  ISS is the appropriate language
because it says precisely that bounded mismatch produces bounded estimation
error, while leaving the magnitude of that error to the anchoring design.

\section{Preliminaries and a Common Error Model}

\subsection{Uniform exponential stability and ISS}
For a discrete LTV homogeneous system
\begin{equation}
  \vct e_{k+1}=\mat A_k\vct e_k,
\end{equation}
let $\mat\Psi(k,j)$ denote its transition matrix.  The origin is uniformly
exponentially stable (UES) if there are constants $M\ge1$ and $0<\rho<1$ such
that
\begin{equation}
  \norm{\mat\Psi(k,j)}\le M\rho^{k-j},\qquad k\ge j.
  \label{eq:ues-def}
\end{equation}
For nonlinear discrete-time systems we use the standard local ISS notion of
Jiang and Wang \cite{JiangWang2001_DiscreteISS}: in a neighborhood of the
origin the state is bounded by a decaying function of the initial condition
plus a class-$\mathcal K$ function of the disturbance supremum.

The EKF-family error recursions considered below can be written locally as
\begin{equation}
 \vct e_{k+1}
 =\mat A_k\vct e_k+\vct\phi_k(\vct e_k)+\mat D_k\vct d_k,
 \qquad
 \norm{\vct\phi_k(\vct e)}\le c_\phi\norm{\vct e}^2,
 \label{eq:local-error-model}
\end{equation}
where $\vct d_k$ collects sensor/model error, imperfect virtual constraints and
bounded exogenous scheduling error.  The quadratic remainder is the normal
local consequence of a continuously differentiable nonlinear measurement and
propagation model after subtracting its first-order linearization.

\begin{lemma}[UES linear part implies a local ISS envelope]
\label{lem:ues-local-iss}
Suppose \eqref{eq:ues-def} holds, $\sup_k\norm{\mat D_k}\le d_D$, and
\eqref{eq:local-error-model} holds on $\norm e\le r$.  If
\begin{equation}
  \frac{M c_\phi r}{1-\rho}<1,
  \label{eq:small-gain-radius}
\end{equation}
then a sufficiently small initial condition and sufficiently small bounded input
remain inside that ball, and the nonlinear recursion is locally ISS.
\end{lemma}

\begin{proof}
Variation of constants gives
\begin{align}
 \norm{e_k}
 &\le M\rho^{k-j}\norm{e_j}
 +\sum_{i=j}^{k-1}M\rho^{k-1-i}
 \bigl(c_\phi\norm{e_i}^2+d_D\norm{d_i}\bigr).
 \label{eq:voc-bound}
\end{align}
Inside $\norm e_i\le r$, replace $\norm e_i^2$ by
$r\norm e_i$.  The induced convolution gain of the UES kernel is at most
$M/(1-\rho)$.  Condition \eqref{eq:small-gain-radius} therefore makes the
state-dependent feedback gain strictly smaller than one.  Standard discrete
small-gain/continuation then yields forward invariance of a smaller ball and a
bound consisting of a decaying initial-condition term plus a finite multiple of
$\sup_i\norm d_i$.  This is local ISS.  This is the same constructive step used
in the companion OU--III proof rather than an appeal to EKF stability by name.
\end{proof}

\subsection{What a zero pseudo-measurement actually means}
Let a virtual measurement be implemented as
\begin{equation}
  \vct y^{0}=0=\mat C\vct x+\vct\nu.
  \label{eq:pseudo-zero}
\end{equation}
If the physical trajectory does not satisfy $\mat C\vct x=0$ instantaneously,
then the correct error model contains the pseudo-measurement mismatch
\begin{equation}
  \vct d^{0}:=-\mat C\vct x_{\rm true}.
\end{equation}
Thus a zero pseudo-measurement is an anchoring model, not a claim that the true
wave state vanishes sample by sample.

\begin{proposition}[Stability--accuracy separation]
\label{prop:stability-accuracy}
If the homogeneous estimator error is UES and the pseudo-measurement mismatch is
bounded, then the resulting estimation error is bounded proportionally to that
mismatch.  Increasing pseudo-measurement information can improve the homogeneous
decay rate while increasing wave-band bias when the zero model is imperfect.
\end{proposition}

\begin{proof}
The mismatch enters \eqref{eq:local-error-model} through the innovation channel
and is therefore one component of $d_k$.  Apply Lemma~\ref{lem:ues-local-iss}.
The stability conclusion depends on boundedness of the input, whereas the
steady error magnitude depends on the closed-loop input gain.  Nothing in ISS
requires that gain to decrease monotonically as the virtual measurement is made
stronger.
\end{proof}

This proposition is the mathematical reason not to infer filter accuracy from
a stability proof.  The OU regularization covariance, PII restoring bandwidth,
and NLO virtual-measurement gain remain performance parameters even after their
stability role has been established.

\section{A General Anchored-Chain Observability Theorem}
\label{sec:chain-theorem}

Consider one scalar kinematic axis with $m$ integrator levels followed by an OU
(or pure-integrator) acceleration state:
\begin{align}
 \dot x_0&=x_1,\quad \dot x_1=x_2,\quad\ldots,\quad
 \dot x_{m-2}=x_{m-1},\\
 \dot x_{m-1}&=a,\qquad \dot a=-\lambda(t)a,
 \qquad 0\le\lambda(t)\le\bar\lambda<\infty.
 \label{eq:general-chain}
\end{align}
The OU case has $\lambda=1/\tau$ with $0<\tau_-\le\tau(t)\le\tau_+$;
the limit $\lambda=0$ is a pure integrated-force state.  Define
\begin{equation}
 E(t)=\exp\!\left(-\int_0^t\lambda(s)\,ds\right),
 \qquad
 A_m(t)=\frac{1}{(m-1)!}\int_0^t(t-s)^{m-1}E(s)\,ds .
 \label{eq:Am-def}
\end{equation}
Then
\begin{equation}
 x_0(t)=\sum_{j=0}^{m-1}\frac{t^j}{j!}x_j(0)+A_m(t)a(0),
 \qquad A_m^{(m)}(t)=E(t)>0.
 \label{eq:head-solution}
\end{equation}

\begin{theorem}[Bounded-gap head anchoring]
\label{thm:head-anchor}
Suppose $x_0$ is observed at $m+1$ instants
$0=t_0<t_1<\cdots<t_m$ whose consecutive gaps obey
$T_-\le t_j-t_{j-1}\le T_+$ with $T_->0$.  Then the
$(m+1)$-state chain \eqref{eq:general-chain} is uniformly observable on that
horizon.  More precisely, the finite-horizon observation matrix has determinant
bounded away from zero uniformly over every admissible $\lambda(\cdot)$ and
sampling pattern.
\end{theorem}

\begin{proof}
Using \eqref{eq:head-solution}, the observation matrix has rows
\begin{equation}
 \left[1,\ t_i,\ \frac{t_i^2}{2!},\ldots,
 \frac{t_i^{m-1}}{(m-1)!},\ A_m(t_i)\right].
 \label{eq:general-observation-matrix}
\end{equation}
Up to the fixed factorial column scalings, its determinant is the Vandermonde
factor
$V(t)=\prod_{i<j}(t_j-t_i)$ multiplied by the $m$th divided difference
$A_m[t_0,\ldots,t_m]$.  The Hermite--Genocchi formula expresses this divided
difference as a simplex average of $A_m^{(m)}=E$.  Since
$0\le t\le t_m\le mT_+$,
\begin{equation}
 E(t)\ge e^{-\bar\lambda mT_+}>0.
\end{equation}
The bounded-gap hypothesis also gives
$\abs{V(t)}\ge T_-^{m(m+1)/2}$.  Hence
\begin{equation}
 \abs{\det\mat O_m}
 \ge c_m T_-^{m(m+1)/2}e^{-\bar\lambda mT_+}>0
 \label{eq:general-det-bound}
\end{equation}
for a constant $c_m>0$ depending only on the fixed factorial scalings and the
simplex mass.  All entries are uniformly bounded on the admissible compact set,
so a positive determinant lower bound gives a positive lower singular-value
bound and therefore a finite-horizon observability-Gramian lower bound.
\end{proof}

\begin{remark}
The theorem explains why increasing the order of the integration chain costs
pseudo-measurement history.  A single head measurement of a four-state
$[S,p,v,a]$ chain requires four distinct instants; direct measurement of an
intermediate state reduces that requirement.  OU--II exploits exactly this by
measuring both $p$ and $v$.
\end{remark}

\begin{lemma}[Stable-tail augmentation preserves detectability]
\label{lem:stable-tail}
Let an LTV system be partitioned into a uniformly observable subsystem and an
unobserved subsystem whose homogeneous transition is UES.  If all coupling
blocks are uniformly bounded, the full system is uniformly detectable.
\end{lemma}

\begin{proof}
Choose a uniformly stabilizing observer injection for the observable block and
zero injection on the stable tail.  The resulting error transition is block
triangular.  Its diagonal blocks are UES, and the off-diagonal response is a
convolution of two exponentially decaying transition bounds with a uniformly
bounded coupling.  That convolution is exponentially bounded, so the complete
observer error is UES.  This is precisely the construction used for the slow
Gauss--Markov accelerometer-bias residual in the OU filters.
\end{proof}

\section{Kalman Pseudo-Measurements: UES to Local ISS}
\label{sec:kalman-general}

The discrete LTV Kalman implication used by the OU papers is the
Anderson--Moore detectability/stabilizability theory, including the
singular-transition extension \cite{AndersonMoore1981_DetectabilityStabilizability,MooreAnderson1980_SingularTransitionStability}.
We state only the specialization required here.

\begin{theorem}[Pseudo-measurement Kalman stability template]
\label{thm:kalman-template}
Consider an admissible Live interval of an error-state filter
\begin{equation}
 \vct e_{k+1}=\mat F_k\vct e_k+\vct w_k,
 \qquad
 \vct y_k=\mat H_k\vct e_k+\vct v_k,
 \label{eq:ltv-kalman}
\end{equation}
where $H_k$ includes the physical and virtual measurement rows.  Assume:
\begin{enumerate}
 \item $F_k$ and $H_k$ are uniformly bounded;
 \item $(F_k,H_k)$ is uniformly detectable;
 \item $(F_k,Q_k^{1/2})$ is uniformly stabilizable;
 \item the process and measurement covariances are uniformly finite and the
       active measurement covariances are uniformly positive definite;
 \item the schedule controlling transition, covariance and pseudo-measurement
       strength is exogenous and bounded, and pseudo-update gaps satisfy the
       finite-horizon conditions used to prove detectability.
\end{enumerate}
Then the stabilizing Riccati/Kalman recursion has a UES homogeneous estimation
error transition.  If the underlying nonlinear filter has a uniformly quadratic
local remainder, the nonlinear Live recursion is locally ISS by
Lemma~\ref{lem:ues-local-iss}.
\end{theorem}

\begin{proof}
The first four hypotheses are the standard LTV Kalman stability conditions in
the cited results and yield a bounded stabilizing Riccati solution and UES
closed-loop homogeneous error transition.  The fifth hypothesis is what makes
the first four uniform rather than pointwise: it prevents arbitrarily long loss
of the anchoring information or unbounded gain/covariance scheduling.  Apply
Lemma~\ref{lem:ues-local-iss} to the nonlinear remainder.
\end{proof}

\subsection{OU--II as a corollary}
For one OU--II axis use state $[v,p,a]^T$ and two joint $v=0,p=0$
pseudo-update instants separated by $T\in[T_-,T_+]$.  Selecting $v,p$ at the
first instant and propagated $v$ at the second gives
\begin{equation}
 \mat O_{\rm II}(T)=
 \begin{bmatrix}
  1&0&0\\0&1&0\\1&0&A_1(T)
 \end{bmatrix},\qquad
 A_1(T)=\int_0^T E(s)\,ds .
 \label{eq:ou2-observation}
\end{equation}
Since $A_1(T)\ge T_-e^{-T_+/\tau_-}>0$, the translational block is UCO.
Two non-collinear inertial/magnetic vector observations over a bounded attitude
transition make the attitude/gyro-bias block observable, while the active
finite-correlation accelerometer-bias residual is a stable tail.

\begin{corollary}[OU--II local ISS]
\label{cor:ou2}
Under the bounded Live-interval assumptions stated in the OU--II companion
analysis---bounded $h$, $\tau$, $\sigma_{aw}$, $r_{p0}$, $r_{v0}$, positive
finite covariances, bounded pseudo-update gaps, non-collinear accepted vector
updates, active finite-$\tau_b$ accelerometer-bias propagation, and no hard
reset---the 18-state OU--II linearized Kalman error is UES and the nonlinear
Live recursion is locally ISS.
\end{corollary}

\begin{proof}
Equation~\eqref{eq:ou2-observation} proves UCO of each translational axis.  The
three axes and attitude/gyro-bias block form a UCO performance subsystem.  The
residual accelerometer-bias block contracts as
$e^{-h/\tau_b}$ and Lemma~\ref{lem:stable-tail} gives full detectability.
Positive finite process driving supplies uniform stabilizability.  Apply
Theorem~\ref{thm:kalman-template}.
\end{proof}

\subsection{OU--III as a corollary}
For one OU--III axis reorder the translational state as $[S,p,v,a]^T$.
This is \eqref{eq:general-chain} with $m=3$, and the $S=0$ pseudo-measurement is
exactly the head observation of Theorem~\ref{thm:head-anchor}.  Four consecutive
bounded-gap pseudo-updates therefore make the axis UCO even when $\tau(t)$
varies measurably inside fixed positive bounds.

\begin{corollary}[OU--III local ISS]
\label{cor:ou3}
Under the admissible Live-interval conditions of the OU--III companion proof,
the 21-state OU--III linearized Kalman error is UES and the nonlinear Live
recursion is locally ISS.
\end{corollary}

\begin{proof}
Theorem~\ref{thm:head-anchor} gives UCO of all three $[S,p,v,a_w]$ axes.  The
accepted inertial/magnetic vector geometry gives the attitude/gyro-bias finite
horizon rank condition.  The finite-$\tau_b$ residual accelerometer-bias block
is UES and hence a detectable stable tail.  Bounded positive process driving
provides uniform stabilizability and the configured covariance clamps provide
finite positive measurement covariance.  Apply
Theorem~\ref{thm:kalman-template} and then Lemma~\ref{lem:ues-local-iss}.
\end{proof}

\begin{remark}[What adaptation is allowed]
Neither corollary requires $\tau$, $\sigma_{aw}$ or pseudo-measurement strength
to be constant.  They may vary arbitrarily within the stated compact set as far
as the deterministic LTV theorem is concerned, provided the resulting matrices
remain bounded and the required measurement gaps remain bounded.  The applied
one-sample scheduling delay additionally makes the active schedule predictable
with respect to the current innovation for the stochastic interpretation.  A
state-dependent tuner that closes an unmodeled fast feedback loop would require
a different proof.
\end{remark}

\section{Transfer to the Two-Frame Lie-Group EKF}
\label{sec:tfg}

The TFG filter carries the same physical OU--III variables
$(R,v,p,S,a_w,b_g,b_a)$ but uses a right-invariant two-frame local error.
Near a consistent state, the conventional MEKF error $e$ and the TFG tangent
error $\xi$ are related by a smooth local coordinate map
\begin{equation}
 \xi=\mat T_k e+\mathcal O(\norm e^2),
 \label{eq:tfg-local-map}
\end{equation}
where, for example, the rotational bridge is the rotated and negated MEKF
increment rather than a bare sign flip.  On a compact attitude/state set away
from the logarithm cut locus, $T_k$ and $T_k^{-1}$ are uniformly bounded.

\begin{theorem}[UES is invariant under a bounded time-varying local coordinate change]
\label{thm:coord-ues}
Let the homogeneous linearized error in coordinates $e_k$ be UES and let
$\xi_k=T_k e_k$ with
\begin{equation}
 \sup_k\norm{T_k}\le\bar T,
 \qquad
 \sup_k\norm{T_k^{-1}}\le\bar T_-.
 \label{eq:T-bounds}
\end{equation}
Then the conjugate linear recursion in $\xi$ is UES.
\end{theorem}

\begin{proof}
If $e_k=\Psi_e(k,j)e_j$, then
\begin{equation}
 \xi_k=T_k\Psi_e(k,j)T_j^{-1}\xi_j.
\end{equation}
Hence
\begin{equation}
 \norm{\Psi_\xi(k,j)}
 \le\bar T\bar T_-M\rho^{k-j},
\end{equation}
which is UES.
\end{proof}

\begin{corollary}[Conditional local ISS of TFG]
\label{cor:tfg}
Assume the TFG process and measurement Jacobians are the first-order
representation, in the two-frame tangent coordinates, of the same physical
Live recursion certified by Corollary~\ref{cor:ou3}; the covariance and noise
models are transformed consistently; and the trajectory remains in a compact
chart satisfying \eqref{eq:T-bounds}.  Then the TFG homogeneous linearized
error is UES and its nonlinear recursion is locally ISS.
\end{corollary}

\begin{proof}
Under the Jacobian-consistency assumption, the first-order TFG recursion is the
bounded coordinate conjugate of the physical OU--III linearization.  Apply
Theorem~\ref{thm:coord-ues}.  The group exponential, residuals and omitted
innovation-dependent invariant-Jacobian terms are smooth near a consistent
state, so their discrepancy from first order is quadratic.  Apply
Lemma~\ref{lem:ues-local-iss}.
\end{proof}

The qualification in Corollary~\ref{cor:tfg} is deliberate.  The implemented
marine process is not claimed to be group affine, and therefore no global
log-linear InEKF theorem is invoked.  What the two-frame construction changes
is the local coordinate geometry and finite retraction, not the physical need
for persistent $S$ anchoring, attitude excitation, bounded covariances and a
valid local chart.  The repository's convention and Jacobian tests provide
implementation evidence for the first-order correspondence; they are not a
substitute for the mathematical hypotheses of the corollary.

\section{PII as Continuous Integral-State Anchoring}
\label{sec:pii}

Ignoring the optional very-slow bias channel for the moment, the vertical PII
core obeys
\begin{align}
 \dot S&=p, & \dot p&=v,\\
 \dot v&=u-3rv-3r^2p-r^3S,
 \label{eq:pii-core}
\end{align}
where $u=a_f-b$ is the filtered acceleration input after any bias correction.
For frozen $r>0$ the characteristic polynomial is
\begin{equation}
 s^3+3rs^2+3r^2s+r^3=(s+r)^3,
 \label{eq:pii-poly}
\end{equation}
so the unforced core is exponentially stable.  This is the deterministic
continuous-feedback analogue of anchoring $S,p,v$ to zero rather than a
sampled Kalman pseudo-measurement.

The adaptive implementation changes $r$ slowly.  Pointwise Hurwitz stability of
$A(r(t))$ alone is not enough for an arbitrarily fast time-varying system, so a
schedule-rate condition is required.

\begin{theorem}[ISS of the slowly scheduled PII core]
\label{thm:pii-scheduled}
Let $0<r_-\le r(t)\le r_+<\infty$ be absolutely continuous and define
$\nu(t)=\dot r(t)/r(t)$.  Let
\begin{equation}
 \vct z=\mat D(r)\begin{bmatrix}S&p&v\end{bmatrix}\T,
 \qquad
 \mat D(r)=\diag(r^2,r,1),
 \end{equation}
so that
\begin{equation}
 \dot z=
 \left[r\mat A_0+\nu\mat J\right]z+\mat D(r)\vct B u,
 \label{eq:pii-scaled}
\end{equation}
with
\begin{equation}
 A_0=\begin{bmatrix}0&1&0\\0&0&1\\-1&-3&-3\end{bmatrix},
 \qquad J=\diag(2,1,0),
 \qquad B=\begin{bmatrix}0\\0\\1\end{bmatrix}.
\end{equation}
Let $P=P^T>0$ solve
\begin{equation}
 A_0^TP+PA_0=-I.
 \end{equation}
If
\begin{equation}
 \bar\nu\,\norm{J^TP+PJ}<r_-,
 \qquad \bar\nu:=\sup_t\abs{\dot r/r},
 \label{eq:pii-rate-condition}
\end{equation}
then the unforced scheduled PII core is UES and the forced core is ISS from
$u$ to $(S,p,v)$.
\end{theorem}

\begin{proof}
The scaling gives \eqref{eq:pii-scaled} by direct differentiation.  With
$V=z^TPz$,
\begin{align}
 \dot V
 &=-r\norm z^2
 +\nu z^T(J^TP+PJ)z
 +2z^TPD(r)Bu\\
 &\le-\alpha_z\norm z^2+c_u\norm z\abs u,
 \label{eq:pii-Vdot}
\end{align}
where
$\alpha_z=r_--\bar\nu\norm{J^TP+PJ}>0$ and $c_u$ is finite because $r$ lies in
a compact positive interval.  Young's inequality gives
\begin{equation}
 \dot V\le-\frac{\alpha_z}{2}\norm z^2
 +\frac{c_u^2}{2\alpha_z}\abs u^2,
\end{equation}
which is an ISS Lyapunov inequality.  Uniform bounds on $D(r)$ and $D(r)^{-1}$
transfer the result to $(S,p,v)$.
\end{proof}

The acceleration prefilter
$\dot a_f=-(1/\tau_a)a_f+(1/\tau_a)a_{\rm meas}$ is itself UES for any
$0<\tau_{a,-}\le\tau_a(t)\le\tau_{a,+}$; no $\dot\tau_a$ bound is needed for
that scalar equation.  Cascading it into Theorem~\ref{thm:pii-scheduled}
preserves ISS.  The implementation's long input and parameter smoothing time
constants are therefore not merely heuristic anti-jitter devices: they also
push the scheduled system toward the explicit slow-variation regime
\eqref{eq:pii-rate-condition}.

\begin{remark}[Optional PII bias channel]
With the optional states
$\dot d=(p-d)/\tau_d$ and $\dot b=k_b d-\lambda_b b$, the bias subsystem is
stable for frozen positive $\tau_d,\lambda_b$ and is driven by $p$, while $b$
re-enters the PII core as an input.  Theorem~\ref{thm:pii-scheduled} therefore
reduces certification of the full optional extension to a standard two-block
ISS small-gain condition.  The present paper does not claim that every arbitrary
manual combination of $(k_b,\lambda_b,\tau_d)$ satisfies that condition; the
repository clamps and very small scheduled $k_b$ make it a separately
verifiable implementation condition rather than an automatic consequence of
the repeated-pole design.
\end{remark}

The semi-implicit embedded update is a consistent discretization of
\eqref{eq:pii-core}.  For a compact $r$ interval satisfying the continuous
margin above there exists $h^*>0$ such that all sufficiently small bounded
steps preserve a discrete UES margin.  The IMU step is orders of
magnitude smaller than the inverse PII pole rate, but a formal maximum-step
certificate for every optional clamp/bias configuration is outside the present
claim.

\section{The Bryne--Fossen--Johansen NLO as an ISS Interconnection}
\label{sec:nlo}

Bryne, Fossen and Johansen augment the vertical translational state with the
integrated height $p'_z=\int p_z\,dt$ and use the virtual measurement $p'_z=0$
for a surface vessel \cite{BryneFossenJohansen2014_TimeVarGain}.  Their
translation error is of the form
\begin{equation}
 \dot{\tilde x}=(A-\vartheta(t)K_\theta C)\tilde x+B\tilde d,
 \label{eq:nlo-x}
\end{equation}
where $\vartheta(t)$ is the time-varying gain multiplier and $\theta$ is a
high-gain scaling parameter.  With
\begin{equation}
 L_\theta=\blkdiag(1,\theta^{-1}I,
 \theta^{-2}I,\theta^{-3}I),
 \qquad \eta=L_\theta\tilde x,
 \label{eq:nlo-scaling}
\end{equation}
their Appendix obtains, for constant $\theta$,
\begin{equation}
 \frac{1}{\theta}\dot\eta
 =(A-\vartheta(t)K_0C)\eta+\rho(t,\tilde\zeta),
 \label{eq:nlo-scaled-error}
\end{equation}
where $\tilde\zeta$ denotes the attitude/gyro-bias error and
$\rho$ is the coupling disturbance.  Choosing $K_0=PC^T$ from their Riccati
equation gives a quadratic Lyapunov function
$U=\theta^{-1}\eta^TP^{-1}\eta$ and an inequality of the form
\begin{equation}
 \dot U
 \le-a_\eta\norm\eta^2
 +b_\eta\theta^{-4}\norm\eta\norm{\tilde\zeta},
 \qquad a_\eta,b_\eta>0.
 \label{eq:nlo-U-iss}
\end{equation}
Their attitude/gyro-bias Lyapunov analysis gives the complementary structure
\begin{equation}
 \dot W
 \le-a_\zeta\norm{\tilde\zeta}^2
 +b_\zeta\theta^{3}\norm{\tilde\zeta}\norm\eta,
 \label{eq:nlo-W-iss}
\end{equation}
inside the semiglobal attitude region used in their theorem.  The positive
power $\theta^3$ is expected: the physical specific-force error is recovered
from the fourth scaled translation coordinate by multiplying by $\theta^3$.

\begin{proposition}[ISS interpretation of the published NLO proof]
\label{prop:nlo-iss}
For fixed admissible $\theta$, inequality \eqref{eq:nlo-U-iss} makes the
translational subsystem ISS with respect to the attitude/gyro-bias error, with a
state-gain coefficient proportional to $\theta^{-4}$.  Inequality
\eqref{eq:nlo-W-iss} makes the attitude subsystem locally/semiglobally ISS with
respect to the scaled translational error, with gain proportional to
$\theta^3$.  Their sufficiently-large-$\theta$ requirement is an interconnection
small-gain requirement in addition to the separate stability of each nominal
subsystem.
\end{proposition}

\begin{proof}
Apply Young's inequality to \eqref{eq:nlo-U-iss}:
\begin{equation}
 \dot U\le-\frac{a_\eta}{2}\norm\eta^2
 +\frac{b_\eta^2}{2a_\eta}\theta^{-8}
  \norm{\tilde\zeta}^2.
 \label{eq:nlo-U-young}
\end{equation}
This is an ISS Lyapunov inequality.  The same step on
\eqref{eq:nlo-W-iss} yields an ISS inequality in the other direction.  The
product of the corresponding state gains scales as $\theta^{-1}$, so increasing
$\theta$ weakens the closed interconnection.  The published composite choice
\begin{equation}
 Y=U+\theta^{-7}W
 \label{eq:nlo-composite}
\end{equation}
is consistent with this interpretation: multiplying the
$\theta^3\norm{\tilde\zeta}\norm\eta$ cross term in $\dot W$ by
$\theta^{-7}$ places it at the same $\theta^{-4}$ order as the cross term in
$\dot U$.  The resulting $2\times2$ quadratic-form condition becomes positive
definite for sufficiently large $\theta$, which is exactly the small-gain
closure used to establish their USGES result.
\end{proof}

This reading is useful because it distinguishes three statements that are easy
to conflate.  First, $A-\vartheta K_0C$ can remain Hurwitz for positive gains.
Second, the translation subsystem can be ISS with respect to attitude error.
Third, the \emph{interconnected} nonlinear observer is USGES only after the
cross-gains satisfy the stronger large-$\theta$ condition and the attitude
initial error lies in the theorem's semiglobal set.  The published theorem
proves the third statement under its assumptions; it is not merely a pole test.

\subsection{What changes when $\theta$ itself is scheduled}
The repository's wave-oriented NLO does not keep the high-gain scale fixed: it
schedules $\theta$ from a wave-frequency proxy.  This is not the same operation
as the published $\vartheta(t)$ multiplier.  If $\theta=\theta(t)$, then
$\dot L_\theta$ in \eqref{eq:nlo-scaling} is no longer zero.  With
\begin{equation}
 \nu_\theta=\frac{\dot\theta}{\theta},
 \qquad
 D_\theta=\blkdiag(0,I,2I,3I),
\end{equation}
one obtains
\begin{equation}
 \dot L_\theta L_\theta^{-1}=-\nu_\theta D_\theta.
 \label{eq:nlo-Ldot}
\end{equation}
Thus the scaled translation dynamics acquire the additional term
$-\nu_\theta D_\theta\eta$.  The Lyapunov function also contains the explicit
factor $1/\theta$, so its derivative receives another bounded term proportional
to $\abs{\nu_\theta}U$.

\begin{theorem}[Conditional local ISS with scheduled NLO high-gain scale]
\label{thm:nlo-scheduled}
Assume $0<\theta_-\le\theta(t)\le\theta_+<\infty$, the published Riccati and
$\vartheta(t)$ lower-bound conditions hold, the attitude variables remain in a
compact local region where the coupling estimates used in
\eqref{eq:nlo-U-iss}--\eqref{eq:nlo-W-iss} are valid, and
$\abs{\dot\theta/\theta}\le\bar\nu_\theta$.  Then there is a constant
$c_\theta>0$ such that the translation Lyapunov derivative satisfies
\begin{equation}
 \dot U
 \le-\bigl(a_\eta-c_\theta\bar\nu_\theta\bigr)\norm\eta^2
 +b_\eta\theta_-^{-4}\norm\eta\norm{\tilde\zeta}.
 \label{eq:nlo-scheduled-bound}
\end{equation}
Consequently, if $c_\theta\bar\nu_\theta<a_\eta$ the translation subsystem
remains ISS.  If, in addition, the two ISS cross-gains satisfy a local small-gain
condition, the scheduled interconnected observer is locally ISS.
\end{theorem}

\begin{proof}
Differentiate $\eta=L_{\theta(t)}\tilde x$ and use
\eqref{eq:nlo-Ldot}.  Every new schedule-rate contribution is bilinear in
$\eta$ and bounded by a constant times
$\abs{\nu_\theta}\norm\eta^2$ on the compact $\theta$ interval.  Differentiating
the prefactor $1/\theta$ in $U$ adds a term with the same bound.  Absorb all of
these terms into $c_\theta\bar\nu_\theta\norm\eta^2$, giving
\eqref{eq:nlo-scheduled-bound}.  A positive residual decay margin yields an ISS
Lyapunov inequality by Young's inequality.  Closing the two subsystems uses the
standard local ISS small-gain argument.
\end{proof}

\begin{remark}[Scope of the repository NLO]
The published theorem assumes a sufficiently large constant high-gain
$\theta\ge\theta^*$; the wave-tuned repository configuration can deliberately
operate below that conservative high-gain range to keep the aiding corner below
the wave band.  Therefore the published USGES theorem cannot simply be quoted
as a certificate for that modified operating policy.  The correct claim is the
conditional local one in Theorem~\ref{thm:nlo-scheduled}: bounded positive
$\theta$, sufficiently slow scheduling, and a verified local interconnection
margin.  This distinction preserves the performance motivation for low
$\theta$ without overstating what the 2014 proof establishes.
\end{remark}

\section{A Unified Stability Map}

\begin{table*}[t]
\centering
\caption{Common anchoring structure and stability route.  ``Virtual state'' means the zero relation is a modeling constraint whose mismatch enters the ISS disturbance.}
\label{tab:unified-map}
\footnotesize
\setlength{\tabcolsep}{4pt}
\begin{tabular}{@{}p{0.13\textwidth}p{0.19\textwidth}p{0.20\textwidth}p{0.20\textwidth}p{0.20\textwidth}@{}}
\toprule
Estimator & Anchored quantity & Linear/core stability mechanism & Nonlinear or scheduled extension & Certified scope in this paper \\
\midrule
OU--II & periodic $p=0$, $v=0$ & UCO of $[p,v,a_w]$; bounded Riccati gives UES & quadratic MEKF remainder & local ISS on admissible Live intervals \\
OU--III & periodic $S=0$ & Theorem~\ref{thm:head-anchor} for $[S,p,v,a_w]$; bounded Riccati gives UES & quadratic MEKF remainder & local ISS on admissible Live intervals \\
TFG & invariant $S=0$ residual & bounded local coordinate conjugacy to OU--III linearization & group retraction and invariant residual remainder & conditional local ISS in a compact chart \\
Vertical PII & continuous restoring of $S,p,v$ & $(s+r)^3$ for fixed $r$; scaled Lyapunov proof for $r(t)$ & slow schedule and stable acceleration LPF & ISS under explicit schedule-rate condition; optional bias loop needs small-gain check \\
Time-varying-gain NLO & virtual integrated-height state $p'_z=0$ & Riccati-designed deterministic observer gain & ISS interconnection of translation and attitude; scheduled $\theta$ adds $\dot\theta/\theta$ term & published constant-$\theta$ observer: USGES under cited assumptions; repository scheduled-$\theta$: conditional local ISS \\
\bottomrule
\end{tabular}
\end{table*}

Table~\ref{tab:unified-map} suggests a reusable design procedure for marine
inertial observers:
\begin{enumerate}
 \item identify the neutral integration modes;
 \item choose a physically defensible virtual quantity whose long-term value is
       bounded or approximately zero;
 \item prove that the resulting output makes the integration chain uniformly
       observable over bounded update gaps;
 \item isolate any unobserved states and require them to be intrinsically stable
       or separately observable;
 \item keep adaptive schedules in compact sets and, when the coordinate/gain
       transformation itself changes, bound its relative rate;
 \item establish UES of the homogeneous linear/core recursion;
 \item treat model error, virtual-constraint mismatch and higher-order nonlinear
       terms as ISS inputs rather than silently setting them to zero.
\end{enumerate}

The order matters.  A pseudo-measurement covariance or restoring gain cannot
repair an unobservable chain if the update disappears for arbitrarily long
periods; conversely, observability alone does not ensure a numerically sensible
or wave-faithful gain.  Stability is a structural prerequisite, not the tuning
objective.

\section{Hybrid Operations and Failure Boundaries}

All proofs above concern continuous Live intervals or their discrete-time
counterparts.  Real implementations contain startup handoffs, covariance
initialization, attitude resets, magnetic re-acquisition, confidence gating,
state clamps and possible measurement rejection.  Those operations are hybrid
jumps and require either separate jump-map inequalities or an interval-by-
interval interpretation.

For the OU filters, the existing local proofs deliberately begin after startup
and re-lock events and require bounded pseudo-update gaps.  TFG inherits the
same restriction.  PII state clamps make the actual implementation globally
bounded in a practical sense when enabled, but saturation is nonlinear and can
invalidate the nominal Lyapunov derivative; boundedness from a clamp should not
be relabeled asymptotic stability.  For NLO, parameter projection bounds the
estimated gyro bias in the published proof, while the scheduled-$\theta$
extension additionally needs the rate condition above.

A long outage is especially revealing.  If the pseudo-measurement is removed,
the low-order inertial chain regains neutral integration modes.  The estimator
may remain bounded for a finite outage because the previous state is finite,
but no theorem that relies on uniform complete observability can cover
arbitrarily long loss of the anchor.  This is a hypothesis failure, not a
contradiction of the theorem.

\section{Discussion}

The main unification is stronger than saying that all five methods ``remove
drift.''  They create a feedback path from an integrated quantity back into an
otherwise neutral kinematic chain.  In OU--II and OU--III the feedback gain is
chosen indirectly by a covariance and Riccati recursion.  In TFG the same
information is expressed in a different local geometry.  In PII it is explicit
pole placement.  In the NLO it is a deterministic high-gain observer driven by
a virtual integrated-position measurement.  The mathematical questions are
therefore the same: Is the chain observable from the anchor?  Is the gain law
uniformly stabilizing?  Is the adaptive schedule bounded and sufficiently
regular?  Are unobserved modes stable?  How large is the mismatch between the
virtual constraint and the physical motion?

The general chain theorem also clarifies the role of OU regularization.  The OU
state is not what makes the integral chain observable.  Its exponential tail
simply replaces the final polynomial basis function by $A_m(t)$ whose $m$th
derivative remains strictly positive.  The same determinant argument includes
the pure-integrator limit.  OU dynamics matter for stochastic prior quality and
regularization tuning, whereas the existence of the finite-horizon anchor is a
separate structural fact.

For TFG, the result is intentionally local.  A Lie-group state can improve
consistency of finite corrections and remove estimated-state dependence from
some first-order measurement Jacobians, but geometry does not abolish the need
for excitation and anchoring.  Because the implemented marine process is not
fully group affine, global log-linear invariant-filter stability is not
available merely from the choice of group.

For PII and NLO, the schedule-rate terms are the most useful practical lesson.
A family of frozen gains can be stable at every operating point while an
arbitrarily fast schedule destabilizes the time-varying system.  The scaled
coordinates expose the exact perturbation: $\dot r/r$ for PII and
$\dot\theta/\theta$ for the scheduled NLO.  The repository's deliberately slow
PII adaptation and bumpless NLO parameter changes therefore have a stability
interpretation in addition to their transient-performance purpose.

Finally, the ISS view gives a clean way to discuss real-ocean deviations from
the zero-mean assumptions.  Tide, current, sustained vessel acceleration,
nonzero mean set-down, imperfect gravity compensation, and slowly varying
sensor bias need not be declared absent.  They can be carried as bounded inputs.
The theorem then says the estimator will not turn a bounded modeling error into
unbounded drift inside the certified regime.  How much of that error leaks into
the reported wave signal is a frequency-response and tuning question that must
be validated separately.

\section{Conclusion}

A common stability structure underlies the OU--II, OU--III, TFG, vertical PII
and time-varying-gain nonlinear marine estimators.  Persistent integral-state
anchoring converts drifting integration chains into observable or directly
stabilized systems; stable unobserved tails preserve detectability; bounded
Riccati recursions or direct Lyapunov feedback then provide UES; and nonlinear
model mismatch is naturally handled through local ISS.  The same framework
also exposes where the methods genuinely differ.  TFG requires a local
coordinate-conjugacy argument rather than a new global invariant theorem.  PII
requires a schedule-rate condition when its repeated pole moves.  The published
NLO theorem proves USGES for a sufficiently large constant high-gain scale and a
time-varying multiplier, whereas scheduling the high-gain scale itself creates
an additional $\dot\theta/\theta$ perturbation and supports only a conditional
local ISS claim unless a new interconnection margin is verified.

The resulting design rule is simple: first prove that the virtual anchor
actually observes the drift chain; then prove that the chosen gain mechanism is
uniformly stabilizing under its real schedule; finally treat the fact that the
virtual zero is only approximately true as an input, not as an identity.  This
separates stability, adaptation and wave fidelity cleanly enough that the same
proof architecture can be reused across estimator families.

\section*{Acknowledgments}
This work was prepared with the supervised assistance of artificial intelligence
tools, including ChatGPT, used for mathematical drafting, source comparison and
editorial preparation.  All conceptual choices, validation and final revisions
remain the responsibility of the author.

\section*{Disclosure Statement}
The author declares no conflicts of interest and received no funding for this
research.

\bibliographystyle{IEEEtran}
\bibliography{w3d-iss,pseudo-meas-stability}

\end{document}
